\documentclass{beamer}
\mode<presentation>
\usepackage{amsmath}
\usepackage{amssymb}
%\usepackage{advdate}
\usepackage{graphicx}
\usepackage{adjustbox}
\usepackage{subcaption}
\usepackage{enumitem}
\usepackage{multicol}
\usepackage{mathtools}
\usepackage{listings}
\usepackage{url}
\def\UrlBreaks{\do\/\do-}
\usetheme{Boadilla}
\usecolortheme{lily}
\setbeamertemplate{footline}
{
  \leavevmode%
  \hbox{%
  \begin{beamercolorbox}[wd=\paperwidth,ht=2.25ex,dp=1ex,right]{author in head/foot}%
    \insertframenumber{} / \inserttotalframenumber\hspace*{2ex} 
  \end{beamercolorbox}}%
  \vskip0pt%
}
\setbeamertemplate{navigation symbols}{}

\providecommand{\nCr}[2]{\,^{#1}C_{#2}} % nCr
\providecommand{\nPr}[2]{\,^{#1}P_{#2}} % nPr
\providecommand{\mbf}{\mathbf}
\providecommand{\pr}[1]{\ensuremath{\Pr\left(#1\right)}}
\providecommand{\qfunc}[1]{\ensuremath{Q\left(#1\right)}}
\providecommand{\sbrak}[1]{\ensuremath{{}\left[#1\right]}}
\providecommand{\lsbrak}[1]{\ensuremath{{}\left[#1\right.}}
\providecommand{\rsbrak}[1]{\ensuremath{{}\left.#1\right]}}
\providecommand{\brak}[1]{\ensuremath{\left(#1\right)}}
\providecommand{\lbrak}[1]{\ensuremath{\left(#1\right.}}
\providecommand{\rbrak}[1]{\ensuremath{\left.#1\right)}}
\providecommand{\cbrak}[1]{\ensuremath{\left\{#1\right\}}}
\providecommand{\lcbrak}[1]{\ensuremath{\left\{#1\right.}}
\providecommand{\rcbrak}[1]{\ensuremath{\left.#1\right\}}}
\theoremstyle{remark}
\newtheorem{rem}{Remark}
\newcommand{\sgn}{\mathop{\mathrm{sgn}}}
\providecommand{\abs}[1]{$\left\vert#1\right\vert$}
\providecommand{\res}[1]{\Res\displaylimits_{#1}} 
\providecommand{\norm}[1]{\lVert#1\rVert}
\providecommand{\mtx}[1]{\mathbf{#1}}
\providecommand{\mean}[1]{E$\left[ #1 \right]$}
\providecommand{\fourier}{\overset{\mathcal{F}}{ \rightleftharpoons}}
%\providecommand{\hilbert}{\overset{\mathcal{H}}{ \rightleftharpoons}}
\providecommand{\system}[1]{\overset{\mathcal{#1}}{ \longleftrightarrow}}
%\providecommand{\system}{\overset{\mathcal{H}}{ \longleftrightarrow}}
	%\newcommand{\solution}[2]{\textbf{Solution:}{#1}}
%\newcommand{\solution}{\noindent \textbf{Solution: }}
\providecommand{\dec}[2]{\ensuremath{\overset{#1}{\underset{#2}{\gtrless}}}}
\newcommand{\myvec}[1]{\ensuremath{\begin{pmatrix}#1\end{pmatrix}}}
\let\vec\mathbf

\lstset{
%language=C,
frame=single, 
breaklines=true,
columns=fullflexible
}

\numberwithin{equation}{section}

\title{12.206}
\author{AI25BTECH11024 - Pratyush Panda}
\begin{document}
\maketitle

\begin{frame}
\textbf{Question: } \\
The eigenvectors of the matrix 
\begin{align}
\myvec{1 & 2 \\ 0 & 2}
\end{align}
are written in the form $\myvec{1 \\ a}$ and $\myvec{1 \\ b}$. What is $a+b$?
\begin{enumerate}
\item[a] 0
\item[b] $\frac{1}{2}$
\item[c] 1
\item[d] 2
\end{enumerate}
\end{frame}

\begin{frame}
\textbf{Solution: } \\
Given
\begin{align}
\Vec{A} = \myvec{1 & 2 \\ 0 & 2}
\end{align}

We have to find the eigen vectors. For that we first have to find the eigen values using the characteristic equation

\begin{align}
det(\Vec{A} - \lambda \Vec{I}) = 0 \\
(1 - \lambda)(2 - \lambda) = 0 \\
\end{align}

Hence,
\begin{align}
\lambda_1 = 1, \quad \lambda_2 = 2
\end{align}
\end{frame}

\begin{frame}
Now, to find the eigenvector for each eigen value;

\begin{align}
\Vec{A}\Vec{v} = \lambda\Vec{v} \\
\implies (\Vec{A} - \lambda I)\vec{v} &= 0 \\
\myvec{1-\lambda & 2 \\ 0 & 2 - \lambda}\Vec{v} = \myvec{0 \\ 0}
\end{align}

Writing as an augmented matrix and row-reducing:
\begin{align}
\myvec{1-\lambda & 2 & \vline & 0 \\ 0 & 2-\lambda & \vline & 0}
\end{align}
\end{frame}

\begin{frame}
For $\lambda=1$ we have;
\begin{align}
\myvec{0 & 2 & \vline & 0 \\ 0 & 1 & \vline & 0}
\end{align}

On solving, we get $\Vec{v}=\myvec{x \\ 0}$. \\
Thus, the first eigen vector is $\myvec{1 \\ 0}$.

For $\lambda=2$ we have;
\begin{align}
\myvec{-1 & 2 & \vline & 0 \\ 0 & 0 & \vline & 0}
\end{align}

On solving this, we get $\Vec{v}=\myvec{x \\ \frac{x}{2}}$. \\
Thus, the second eigen vector is $\myvec{1 \\ \frac{1}{2}}$.
\end{frame}

\begin{frame}
Thus, when we get $a=0$ and $b=\frac{1}{2}$. So $a+b=\frac{1}{2}$.
Therefore, option d is correct.
\end{frame}

\end{document}